\par\noindent\SourceLeaf{51}%
\hspace*{\parindent}6. Nunc sine mora ad finem propositum accedere licet. Nimirum evolvatur per theorema binomiale $(1+n\cos\varphi)^{\nu}$ in seriem per potestates ipsius $\cos\varphi$ progredientem, solvatur deinde quodlibet huius seriei membrum in novam seriem secundum cosinus multiplorum anguli dati procedentem, seligantur ex his seriebus termini in eiusdem anguli cosinum ducti, hique in unum terminum cogantur. Cognita terminorum sic provenientium lege totum negotium tantum non perfectum erit. In serie fundamentali

\[
\begin{aligned}
(1+n\cos\varphi)^{\nu} & = 1 + {}^{\nu}\mathfrak{B}^{(1} n\cos\varphi + {}^{\nu}\mathfrak{B}^{(2)} n^{2}\cos\varphi^{2} \\&\quad + {}^{\nu}\mathfrak{B}^{(3)} n^{3}\cos\varphi^{3} + .\,. + {}^{\nu}\mathfrak{B}^{(q)} n^{q}\cos\varphi^{q} + .\,.
\end{aligned}
\]

\noindent omnes ipsius $\cos\varphi$ potestates ex ordine occurrunt, pares et impares. Ex paribus, si evolvantur, nonnisi cosinus multiplorum parium, ex imparibus nonnisi imparium prodeunt, quare ad has, si de cosinu multipli imparis, ad illas si de cosinu multipli paris agitur, recurrendum est.

A) Habebitur seriei quaesitae terminus initialis, si ex omnibus propositae terminis, illorum, qui evoluti a $\cos 0\varphi = 1$ ordiuntur, initialia membra ordine suo servato uniantur. Sunt vero solae potestates ipsius $\cos\varphi$ pares, quae evolutae eiusmodi membrum initiale praebeant. Quaelibet harum partem huc pertinentem praebet; $h^{\text{ta}}$ igitur pars ex seriei fundamentalis termino $2h^{\text{to}}$, qui est ${}^{\nu}\mathfrak{B}^{(2h)} n^{2h}\cos\varphi^{2h}$, proficiscitur. Est autem $\cos\varphi^{2h}$ terminus initialis, qui solus huc pertinet

\[
= \frac{1}{2^{2h-1}} \frac{{}^{2h}\mathfrak{B}^{(h)}}{2} = \frac{1}{2^{2h}} {}^{2h}\mathfrak{B}^{(h)};
\]

\noindent adeoque adiecto ipsius $\cos\varphi^{2h}$ coefficiente pars $h^{\text{ta}}$ seriei in $\cos 0\varphi = 1$ ductae

\[
\frac{1}{2^{2h}} {}^{2h}\mathfrak{B}^{(h)}  .  {}^{\nu}\mathfrak{B}^{(2h)} n^{2h}.
\]

\noindent Sic inventus est terminus generalis $h^{\text{tus}}$ seriei, quae terminum initialem eius constituit, quam $(1+n\cos\varphi)^{\nu}$ secundum cosinus angulorum multiplorum evolvendo adipiscimur. Esse terminum eius initialem $=1$ per se patet. Quod si ponamus esse

\[
(1+n\cos\varphi)^{\nu} = A + A^{(1)}\cos\varphi + A^{(2)}\cos 2\varphi + .\,.\,.,
\]

\noindent erit

\[
A = 1 + \frac{1}{2^{2}} {}^{2}\mathfrak{B}^{(1)}  .  {}^{\nu}\mathfrak{B}^{(2)} n^{2} + \frac{1}{2^{4}} {}^{4}\mathfrak{B}^{(2)}  .  {}^{\nu}\mathfrak{B}^{(4)} n^{4} + .\,.\,. + \frac{1}{2^{2h}} {}^{2h}\mathfrak{B}^{(h)}  .  {}^{\nu}\mathfrak{B}^{(2h)} n^{2h} + .\,.\,.
\]

\noindent sive

\[
A = 1 + \frac{1}{2^{2}}  .  \frac{2}{1} \frac{\nu.\nu-1}{1.2} n^{2} + \frac{1}{2^{4}}  .  \frac{4.3}{1.2}  .  \frac{\nu.\nu-1.\nu-2.\nu-3}{1.2.3.4} n^{4} + \text{etc.}
\]
